\documentclass[prb,twocolumn,nofootinbib,superscriptaddress]{revtex4-2}

% ---------- Packages ----------
\usepackage{amsmath,amssymb}
\usepackage{graphicx}
\usepackage{xcolor}
\usepackage[colorlinks=true,linktoc=page,linkcolor=purple,citecolor=teal,urlcolor=magenta]{hyperref}

% ---------- Macros ----------
\newcommand{\R}{\mathbb{R}}
\newcommand{\Rp}{\mathbb{R}_{>0}}
\newcommand{\Z}{\mathbb{Z}}
\newcommand{\J}{J}
\newcommand{\phig}{\varphi}
\newcommand{\lean}[1]{\texttt{\small #1}}

\begin{document}

\title{Coherent Energy Concentration in Periodic Lattices:\\
Exact $N$-Factor Amplification from Ratio-Energy Minimization}

\author{Jonathan Washburn}
\email{jon@recognitionphysics.org}
\affiliation{Recognition Physics Research Institute, Austin, Texas 78701, USA}

\begin{abstract}
We prove that periodic configurations on finite connected unit-cell lattices are the unique global minimizers of the reciprocal ratio energy $\J(x)=\frac12(x+x^{-1})-1$, and therefore have exactly zero edge-cost density.
Any nonperiodic configuration must contain at least one edge with strictly positive cost.
Combining this with the compact-phase quotient construction, a perfect crystal corresponds to a phase-uniform state.
We then analyze a linear coherent-observable model in which each unit cell contributes the same amplitude $A_0$.
Within that model, coherent summation gives energy $N^2 A_0^2$, while a random-phase baseline gives $N A_0^2$.
The resulting concentration factor is therefore exactly $N$, is monotone in crystal size, and is bounded below by $8$ on any lattice with at least two cells per axis.
All results are proved as theorems in the Lean~4 interactive theorem prover with zero remaining \texttt{sorry} markers.
The framework yields a falsifiable prediction: under matched normalization, coherent enhancement should track coherent-domain size and decrease with defect density or mosaic breakup, testable via diffraction or Brillouin-style spectroscopy.
\end{abstract}

\maketitle

%=====================================================================
\section{Introduction}
\label{sec:intro}
%=====================================================================

We study the reciprocal ratio energy
\begin{equation}
C_G[x] = \sum_{\langle v,w \rangle \in E} \J\!\left(\frac{x_v}{x_w}\right),
\qquad \J(x) = \frac12\!\left(x + x^{-1}\right) - 1,
\label{eq:energy}
\end{equation}
on positive fields $x: V \to \Rp$ over finite connected graphs $G = (V,E)$.
Two structural features of this framework drive the present paper:
\begin{enumerate}
\item zero-temperature minimizers of~\eqref{eq:energy} are constant on connected lattices;
\item after imposing a discrete scaling quotient $x \sim b^n x$, the induced compact phase $\Theta = \log_b x \bmod 1 \in \R/\Z$ is likewise uniform in the minimized state.
\end{enumerate}
These two facts mean that the ratio-energy framework enforces global uniformity of both the scale field and the induced compact phase on connected lattices.

The present paper asks the physical question that follows immediately from this uniformity: \emph{what quantitative consequence does periodicity have for coherent energy concentration in a crystal lattice?}

The new content is fourfold.
First, we state the crystal-specific corollary of rigidity in a unit-cell lattice model: perfect periodicity gives zero edge-cost density, while any departure from that state introduces at least one positive-cost edge.
Second, we combine the zero-cost crystal with the compact-phase quotient model to obtain a phase-uniform crystal state.
Third, we introduce an explicit linear observation model that compares coherent summation against a random-phase benchmark and proves an exact concentration factor of $N$ within that model.
Fourth, we identify a direct experimental discriminator: the coherent enhancement should track coherent-domain size as defect density increases.

A \emph{crystal} in the present paper is modeled by a scalar unit-cell field in which every unit cell is identical to its neighbors: a periodic repetition of the same local state across the lattice.
By zero-temperature rigidity, such a configuration places every edge ratio at $x_v/x_w = 1$, achieving the unique global minimum $\J(1) = 0$ on every edge.
The crystal therefore has \emph{exactly zero local cost density}---a property that no amorphous (nonperiodic) structure can share.

This zero-cost property has a direct physical consequence once one adopts a linear coherent-observable model: because all unit cells in a crystal contribute to the measured observable at the same phase, amplitudes sum coherently.
An observable contributed by $N$ identical cells in phase has amplitude $N A_0$ and energy $N^2 A_0^2$, while $N$ cells with random phases have an expected energy of only $N A_0^2$.
The ratio---the \emph{concentration factor}---is exactly $N$.
For a macroscopic crystal with $N \sim 10^{23}$ unit cells, this is a factor of $10^{23}$: an enormous coherent amplification that distinguishes crystalline from amorphous matter in the ratio-energy framework.

All results in this paper are stated as theorems and machine-verified in Lean~4 with zero remaining \texttt{sorry} markers.
We also derive a falsifiable experimental prediction (Sec.~\ref{sec:prediction}).

%=====================================================================
\section{Zero cost density of periodic configurations}
\label{sec:zero_cost}
%=====================================================================

We begin with the result that distinguishes crystals from amorphous structures in the ratio-energy framework.

\subsection{Definitions}

A \emph{periodic configuration} in the present scalar unit-cell reduction is a positive field $x: V \to \Rp$ on the graph of unit cells such that $x_v = x_w$ for every edge $\langle v,w \rangle \in E$.
On a connected graph, this is equivalent to $x_v \equiv c$ for some constant $c > 0$.
Physically, this represents the reduction in which each unit cell carries the same scalar local state.
It is not, by itself, a full theorem about arbitrary crystallographic motifs or full structure factors.

An \emph{amorphous configuration} is any positive field that is not periodic.
On a connected graph, this means that there exists at least one edge $\langle v,w \rangle$ with $x_v \neq x_w$.

\subsection{Results}

\vspace{2mm}
\underline{\bf Theorem 1 (Zero cost density).}
\emph{In a periodic configuration, the edge cost vanishes on every edge:}
\begin{equation}
x_v = x_w \quad \Longrightarrow \quad \J(x_v / x_w) = \J(1) = 0.
\label{eq:zero_cost}
\end{equation}
\emph{Lean reference:} \lean{periodic\_Jcost\_zero}.
\vspace{2mm}

This is an immediate consequence of $\J(1) = 0$, which follows from $\J(x) = (x-1)^2/(2x) = 0$ at $x = 1$.

\vspace{2mm}
\underline{\bf Theorem 2 (Positive amorphous cost).}
\emph{If $x > 0$ and $x \neq 1$, then $\J(x) > 0$.}
\begin{equation}
x > 0, \quad x \neq 1 \quad \Longrightarrow \quad \J(x) > 0.
\label{eq:positive_cost}
\end{equation}
\emph{Lean reference:} \lean{amorphous\_cell\_cost\_positive}.
\vspace{2mm}

This follows from $\J(x) = (x-1)^2/(2x)$, which is strictly positive when $x \neq 1$ and $x > 0$.

\vspace{2mm}
\underline{\bf Theorem 3 (Crystal advantage).}
\emph{A periodic configuration has strictly lower cost than any nonperiodic perturbation on every affected edge:}
\begin{equation}
x > 0, \quad x \neq 1 \quad \Longrightarrow \quad \J(1) < \J(x).
\label{eq:advantage}
\end{equation}
\emph{Lean reference:} \lean{crystal\_Jcost\_advantage}.
\vspace{2mm}

Theorems~1--3 together establish that among all positive fields on a connected graph, the periodic (constant) configurations are the unique minimizers of the local cost density, and any departure from periodicity incurs a strict cost penalty.

\paragraph{GCIC corollary.}
The zero-cost crystal acquires a compact-phase interpretation when combined with the quotient model.
There, zero edge cost together with connectedness implies uniform compact phase in the discrete quotient model.
Applying that result to the crystal field produced by Theorem~1 yields a phase-uniform crystal: in the quotient description, all unit cells share the same compact phase $\Theta$.
In the broader Recognition Science language, this is the precise sense in which a macroscopic crystal acts as a static phase-uniform, or GCIC-aligned, body.

%=====================================================================
\section{Coherent energy concentration}
\label{sec:concentration}
%=====================================================================

We now quantify the physical consequence of zero cost density within an explicit observation model.

\subsection{Crystal lattice model}

We model a three-dimensional crystal as a rectangular parallelepiped of unit cells with dimensions $n_1 \times n_2 \times n_3$, where $n_i \geq 2$.
The total number of unit cells is
\begin{equation}
N = n_1 \cdot n_2 \cdot n_3 \geq 8.
\label{eq:N}
\end{equation}
This model is formalized as the structure \lean{CrystalLattice} with the constraint \lean{totalCells\_ge\_8}.

\subsection{Observation model: coherent and incoherent benchmarks}

Each unit cell contributes to a collective observable with amplitude $A_0 > 0$.
In the Lean-certified RS instantiation, the choice is $A_0 = \phig^{-5}$, where $\phig = (1+\sqrt{5})/2$ is the golden ratio.
For the scaling results below, however, only positivity of $A_0$ matters.
The exact $N$-factor result is therefore best understood as an exact consequence of combining the crystal-alignment theorem with a linear coherent-observable model.

\paragraph{Coherent case (crystal).}
Because all unit cells are periodic (identical local state), and because the quotient model assigns them a common compact phase, they contribute at the same phase.
The total amplitude is $N A_0$ and the total energy is
\begin{equation}
E_{\mathrm{coh}} = (N A_0)^2 = N^2 A_0^2.
\label{eq:E_coh}
\end{equation}

\paragraph{Incoherent benchmark.}
As a comparison baseline, consider a random-phase superposition in which the phases of the $N$ unit-cell contributions are uncorrelated.
Then the expected energy is
\begin{equation}
E_{\mathrm{inc}} = N A_0^2.
\label{eq:E_inc}
\end{equation}

\subsection{Concentration factor}

Within this observation model, the ratio of coherent to incoherent energy is the \emph{concentration factor}:
\begin{equation}
\eta(N) \equiv \frac{E_{\mathrm{coh}}}{E_{\mathrm{inc}}} = \frac{N^2 A_0^2}{N A_0^2} = N.
\label{eq:eta}
\end{equation}

\vspace{2mm}
\underline{\bf Theorem 4 (Concentration factor is exactly $N$).}
\emph{Within the coherent/incoherent benchmark model above, for any crystal lattice with $N$ unit cells,}
\begin{equation}
\eta(N) = N.
\end{equation}
\emph{Lean reference:} \lean{concentrationFactor\_eq\_totalCells}.
\vspace{2mm}

\vspace{2mm}
\underline{\bf Theorem 5 (Crystal focuses energy).}
\emph{Within the same benchmark model, for any crystal lattice with $N \geq 8$, the coherent energy strictly exceeds the incoherent energy:}
\begin{equation}
E_{\mathrm{inc}} < E_{\mathrm{coh}}.
\end{equation}
\emph{Lean reference:} \lean{crystal\_focuses\_theta\_energy}.
\vspace{2mm}

\vspace{2mm}
\underline{\bf Theorem 6 (Minimum concentration).}
\emph{The concentration factor is bounded below by $8$ for any crystal lattice:}
\begin{equation}
\eta(N) \geq 8.
\end{equation}
\emph{Lean reference:} \lean{concentrationFactor\_ge\_8}.
\vspace{2mm}

\vspace{2mm}
\underline{\bf Theorem 7 (Monotonicity).}
\emph{The concentration factor is monotone increasing in crystal size:}
\begin{equation}
N_1 \leq N_2 \quad \Longrightarrow \quad \eta(N_1) \leq \eta(N_2).
\end{equation}
\emph{Lean reference:} \lean{concentration\_monotone}.
\vspace{2mm}

%=====================================================================
\section{Collective cost subadditivity}
\label{sec:collective}
%=====================================================================

The concentration factor describes how a crystal amplifies a coherent signal relative to an incoherent baseline.
A complementary result concerns the cost comparison between aligned and misaligned collections of boundaries.

\vspace{2mm}
\underline{\bf Corollary 8 (Collective subadditivity).}
\emph{If $N \geq 1$ boundaries have costs $c_1, \ldots, c_N > 0$ (non-unit ratios), then:}
\begin{equation}
\J(1) = 0 < \sum_{i=1}^{N} c_i.
\end{equation}
\emph{That is, the perfectly aligned configuration (all ratios at unity, collective cost zero) strictly beats any collection of misaligned boundaries.}

\emph{Lean reference:} \lean{aligned\_beats\_any\_perturbation}.
\vspace{2mm}

This result is algebraically elementary, but conceptually useful: it replaces what had previously been an empirical hypothesis about collective amplification.
It follows directly from two proved properties of $\J$: (i) $\J(1) = 0$ and (ii) $\J(x) > 0$ for $x \neq 1$, $x > 0$.
No additional axiom or empirical input is required.

%=====================================================================
\section{Machine verification}
\label{sec:lean}
%=====================================================================

All theorems in this paper have been formalized and verified in Lean~4 using Mathlib, with zero remaining \texttt{sorry} markers.
The formalization is organized as follows:

\begin{itemize}
\item \lean{GraphRigidity}: ratio rigidity, edge cost characterization.
\item \lean{ReducedPhasePotential}: reduced phase potential $\widetilde{\J}_b$, periodicity, stiffness.
\item \lean{GCICDerivation}: phase uniformity from rigidity and gauge, collective subadditivity.
\item \lean{CrystalThetaFocusing}: crystal lattice, zero cost density, coherent/incoherent energy, concentration factor, monotonicity.
\end{itemize}

The master certificate \lean{CrystalThetaFocusingFullCert} bundles all twelve theorem fields (Theorems~1--8 and supporting lemmas) into a single Lean structure with default proofs for every field, confirming that the entire derivation chain compiles without axioms, hypotheses, or incomplete proofs.

%=====================================================================
\section{Experimental prediction}
\label{sec:prediction}
%=====================================================================

The concentration factor $\eta = N$ yields a sharp, falsifiable prediction.

\paragraph{Prediction.}
A useful experimental discriminator is the dependence of coherent enhancement on coherent-domain size.
Consider a set of samples with the same composition and total mass but different defect density, grain size, or mosaicity.
If one measures a normalized ratio of coherent spectral weight to diffuse background, the present framework predicts that the ratio should track the number of unit cells in a coherent domain.

For a 1\,kg quartz crystal with lattice parameter $a \approx 5$\,\AA, the unit-cell count is of order
\begin{equation}
N \sim \frac{m}{\rho\, a^3}
  \sim \frac{1\,\text{kg}}{2650\,\frac{\text{kg}}{\text{m}^3}
       \cdot (5{\times}10^{-10}\,\text{m})^3}
  \sim 3 \times 10^{24}.
\end{equation}
The ideal coherent-domain concentration factor is therefore $\eta \sim 3 \times 10^{24}$ for a kilogram-scale single crystal.
This exhibits the familiar $N^2$ versus $N$ contrast between coherent and diffuse contributions.
Within the present framework, that contrast is derived from ratio-energy minimization and departures from ideal scaling are interpreted as signatures of defect-induced positive-cost edges.

The $N$-factor amplification is therefore not merely an empirical observation about constructive interference, but a consequence of the cost functional~\eqref{eq:energy} and the uniqueness of its global minimizer within the present model.

\paragraph{Stronger test.}
A more discriminating experiment would compare the ratio of coherent to incoherent spectral weight as a function of crystal quality (defect density, grain size, mosaicity), with normalization chosen so that total sample mass does not trivially dominate the comparison.
The prediction is that the concentration factor decreases from $N$ toward $N_{\text{dom}}$ (the number of cells per coherent domain) as defects break the perfect periodicity, since defects introduce edges with $\J(x_v/x_w) > 0$ and thereby fragment the zero-cost manifold into smaller coherent regions.

%=====================================================================
\section{Discussion}
\label{sec:discussion}
%=====================================================================

The results of this paper establish three things.

First, periodic configurations are the unique zero-cost-density states of the ratio energy.
This is a structural theorem about the functional~\eqref{eq:energy} that does not depend on any physical interpretation: on any finite connected graph, the constant fields are the only configurations for which every edge contributes zero cost.

Second, within the linear coherent-observable model introduced in Sec.~\ref{sec:concentration}, the physical consequence of zero cost density is coherent amplification by a factor of $N$.
This connects the rigidity statement to a concrete, measurable quantity while keeping clear which part of the argument is rigidity-based and which part comes from the observation model.
Crystalline structures are, in this precise sense, \emph{macroscopic coherence amplifiers} within the ratio-energy framework.

Third, collective cost subadditivity---the fact that the aligned collective strictly beats any sum of misaligned individuals---follows from the cost functional alone, with no additional empirical input.
Although elementary, this corollary is useful because it gives a direct theorem-level statement of collective advantage within the same formal framework.

In the quotient model, the crystal's zero cost density means that all $N$ unit cells share the same compact phase $\Theta$.
A macroscopic crystal is therefore a structure in which the phase field is perfectly uniform across $\sim 10^{23}$ sites, amplifying any phase-coupled observable by the full number of unit cells.
Whether this amplification has consequences beyond ordinary diffraction and coherent vibrational measurements---for instance, whether it enables detectable coupling between macroscopic crystals through the phase-rigidity mechanism---is an open experimental question that the present mathematical framework motivates but does not resolve.
The more immediate empirical target is the predicted dependence on coherent-domain size and defect-controlled fragmentation of the zero-cost manifold.

Our analysis is limited to the zero-temperature, static setting used here.
Finite-temperature effects, dynamics, topological defects, and boundary conditions all represent directions for future work within the broader uniformly convex gradient-model setting to which the ratio energy belongs.

%=====================================================================
\begin{acknowledgments}
We thank Sebastian Pardo-Guerra for discussions.
This work utilized the computing resources of the Recognition Physics Research Institute.
We acknowledge the use of AI-assisted tools for drafting and language refinement.
All scientific content, interpretations, and conclusions were reviewed and validated by the authors.
\end{acknowledgments}

\end{document}
